%Questo è il preambolo, dove si inseriscono i pacchetti e le impostazioni che servono per compilare il documento. Quanto scritto dopo il simbolo '%' è solo un commento e serve a fini dimostrativi.
\documentclass[a4paper, 13pt, oneside]{report}
\linespread{1.5} %interlinea
\pagestyle{plain}
\usepackage{geometry} %margini
\geometry{a4paper, top=3cm, bottom=3cm, left=3cm, right=3cm, bindingoffset=5mm}
\usepackage{graphicx}
\graphicspath{immagini/}
\usepackage{multicol} %più colonne
\usepackage{ragged2e} %allineamento testo


\usepackage[italian]{babel} %lingua principale

\usepackage{minitoc} %mini sommario a inizio capitolo
\nomtcrule 
\addto{\captionsitalian}{% Making babel aware of special titles
  \renewcommand{\mtctitle}{Sommario}
}




\usepackage[sorting=none,backend=biber,refsection=chapter]{biblatex} %bibliografia
\defbibheading{bibbysubsect}{\section*{Capitolo \ref{blx:refsection\therefsection}}}
\addbibresource{bibliografia.bib}





\usepackage{comment}








\usepackage{xpatch}
\usepackage{blindtext}


\makeatletter

\xpatchcmd{\@makeschapterhead}{%
  \Huge \bfseries  #1\par\nobreak%
}{%
  \Huge \bfseries\centering #1\par\nobreak%
}{\typeout{Patched makeschapterhead}}{\typeout{patching of @makeschapterhead failed}}


\xpatchcmd{\@makechapterhead}{%
  \huge\bfseries \@chapapp\space \thechapter
}{%
  \huge\bfseries\centering \@chapapp\space \thechapter
}{\typeout{Patched @makechapterhead}}{\typeout{Patching of @makechapterhead failed}}

\makeatother

\usepackage{fancyhdr}
\usepackage[export]{adjustbox}



\usepackage{hyperref} %hyperlink
\hypersetup{
    colorlinks=true,
    citecolor=cyan,
    linkcolor=black,
    urlcolor=black,
    pdftitle=Una buona tesi,
    pdfauthor=Un povero studente di medicina,
    }

\usepackage{booktabs} %per le tabelle
\usepackage{multirow}
\usepackage[table,xcdraw]{xcolor}
\usepackage{graphicx}

\usepackage{lscape}
\usepackage{float}
\usepackage{wrapfig}

%Da qui in poi inizia il documento.
\begin{document}
\begin{titlepage} %Cambia i valori in \vspace{} per ottenere un risultato perfetto

\begin{center}
	{\scshape\Large\bfseries UNIVERSITÀ DEGLI STUDI DI NAPOLI\\ FEDERICO II \par}
	\vspace{1cm}
	\includegraphics{immagini/logo.jpg}\par
	\vspace{0,7cm}
	{\scshape\Large\bfseries SCUOLA DI MEDICINA E CHIRURGIA \par}
	\vspace{0,5cm}
	{\scshape\normalsize Corso di Laurea Magistrale in Medicina e Chirurgia\par}
	\vspace{0.3cm}
	{\scshape\small Dipartimento di Neuroscienze e Scienze riproduttive\\ ed odontostomatologiche\par}
	\vspace{0.3cm}
	{\scshape\normalsize Tesi di Laurea Sperimentale in\\
    PSICHIATRIA\par}
	\vspace{0.3cm}
	{\scshape\normalsize\bfseries INSERISCI QUI IL TITOLO DELLA TUA TESI\par}
	\vspace{0,5cm}
	{\schape\normalsize\bfseries ADD HERE YOUR TESIS TITLE, THIS TIME IN ENGLISH\par}
	\vspace{1,2cm}
	\end{center}
	
	\begin{multicols}{2}
	\begin{center}
	 	{
		\schape\normalsize\textbf{Relatore:\\Ch.mo.}\\ Giovanni Rossi \par}
	\vfill
	\schape\normalsize\textbf{Candidato:}\\Giuseeppe Esposito\\M39000000 \par 
	\end{center}
	\end{multicols}
	
	\vfill

% Bottom of the page
	{\begin{center}
	     \schape\large ANNO ACCADEMICO 2019/2020
	\end{center}}

\end{titlepage}

\dominitoc
\tableofcontents

\begin{center}
\vfill
    \chapter{Piccoli suggerimenti}
    \label{blx:refsection\therefsection}
\vfill
\minitoc
\newpage
\end{center}
\justify

\section{Introduzione}
Ottima scelta! \LaTeX è di gran lunga meglio di MS Word e programmi affini. Per qualsiasi cosa, fai riferimento alla documentazione di Overleaf o cerca con il tuo motore di ricerca preferito (in inglese) quello che ti serve. Qui ti lascio solo qualche consiglio.

\section{Tabelle}
Usa questo sito https://www.tablesgenerator.com/.

\section{Bibliografia}
Ecco una dimostrazione di come fare una citazione \cite{lysaker2010metacognition}. Ricorda che ci sono gli hyperlink, quindi prova a cliccare sulla citazione sul documento compilato. \\
Il modo più semplice per aggiungere articolari da citare al file 'bibliografia.bib' è andare su google scholar, cercare l'articolo che ti interessa, cliccare sul simbolo ' " ' ed infine selezionare BibTex e fare copia-incolla.

\section{Presentazione}
Prova ad usare \LaTeX anche per la presentazione della tua tesi, non te ne pentirai!

\medskip


\printbibheading[
heading=bibintoc,
title={Bibliografia}
]
\bibbysection[heading=bibbysubsect]


\end{document}
